\section{The Decision Quantity, Audited}
\label{sec:statistics}

Everything in this section is computed from the stored per-item outcome vectors
of the same runs reported above, by a single script
(\texttt{analysis/gen\_extended\_stats.py}); no value is transcribed by hand.
The section exists because a decision rule built on a $\max$ invites two
objections that the conference version had no room to answer, and because a
study reporting a null owes the reader its per-group statistics rather than the
two cells that happen to be significant.

\subsection{The pre-registered question: does GRPO clear $+3$?}
\label{subsec:gotest}

The decision rule this study registered is a Go threshold, not a test of
inferiority: GRPO passes if $\Delta \geq +3$ points. That is the hypothesis the
analysis owes the reader an answer to, and it is a weaker and more answerable
one than ``is GRPO significantly worse.'' Conflating the two understates the
result, so we report both, starting with the registered one.

A paired item bootstrap ($B = 50{,}000$, $\max$ recomputed on each resample)
gives, per group, the probability that $\Delta$ reaches the Go threshold:

\begin{center}
\begin{tabular}{lrr}
\toprule
question & criterion & groups \\
\midrule
Does GRPO clear the registered $+3$ threshold? & 95\% CI upper bound $< +3$ & 14/17 \\
Is GRPO strictly worse than $\max\mathcal{B}$?  & 95\% CI upper bound $< 0$  & 5/17 \\
\bottomrule
\end{tabular}
\end{center}

$P(\Delta \geq +3)$ has a median of $0.0013$ across the 17 groups and does not
exceed $0.0804$ in any of them. The three groups whose intervals fail to
exclude $+3$ are all in Config~1---$m{=}2$ seeds 0 and 1, and $m{=}64$ seed
2---where $n{=}282$ and GRPO and the solver disagree on 70 to 82 items, giving a
wide paired interval. Config~3 has the same $n$ but excludes $+3$ in all six
groups, because GRPO and continued SFT agree on all but 21 to 35 items and the
paired interval is correspondingly narrow. Counts at the $+3$ boundary move by
one group between bootstrap runs; the continuous $P(\Delta \geq +3)$ does not,
and should be preferred.

So the registered claim---GRPO does not reach the threshold that would have
made this a positive result---is supported in 14 of 17 groups, and the
remaining three are inconclusive rather than contrary. The stricter claim that
GRPO is \emph{worse} is supported in 5. We make the first claim, not the
second.

\subsection{Per-group significance, and what it does not show}

% Generated by analysis/gen_extended_stats.py. Do not edit by hand.
\begin{table}[t]
\centering
\small
\caption{Per-group significance and power. $b/c$ are the McNemar discordant counts (GRPO right / $\arg\max\mathcal{B}$ right). $p_{\text{Holm}}$ corrects across all 17 groups. The CI is a paired item bootstrap ($B{=}10{,}000$) of $\Delta$ with the $\max$ recomputed on each resample. MDE is the smallest $|\Delta|$ this group's discordant count could have declared significant, uncorrected and under Holm.}
\label{tab:pergroup}
\begin{tabular}{rrlrrrrcrr}
\toprule
$m$ & seed & $\arg\max\mathcal{B}$ & $\Delta$ & $b/c$ & $p$ & $p_{\text{Holm}}$ & 95\% CI & MDE & MDE$_{\text{H}}$ \\
\midrule
\multicolumn{10}{@{}l}{\textit{Config~1 (CodeTAT-QA $\times$ 0.5B, $n{=}282$)}} \\
2 & 0 & solver & $-2.13$ & 38/44 & 0.581 & 1.000 & $[-8.51,+3.90]$ & 7.09 & 9.93 \\
2 & 1 & solver & $-1.42$ & 37/41 & 0.734 & 1.000 & $[-7.45,+4.61]$ & 7.09 & 9.93 \\
2 & 2 & solver & $-3.90$ & 32/43 & 0.248 & 1.000 & $[-9.93,+2.13]$ & 6.74 & 9.57 \\
64 & 0 & solver & $-6.03$ & 30/47 & 0.068 & 0.878 & $[-12.06,+0.00]$ & 6.74 & 9.57 \\
64 & 1 & solver & $-3.55$ & 30/40 & 0.282 & 1.000 & $[-9.22,+2.13]$ & 6.38 & 9.22 \\
64 & 2 & solver & $-2.84$ & 32/40 & 0.410 & 1.000 & $[-8.87,+3.19]$ & 6.38 & 9.93 \\
\addlinespace
\multicolumn{10}{@{}l}{\textit{Config~2 (CodeFinQA $\times$ 1B, $n{=}664$)}} \\
128 & 0 & solver & $-15.36$ & 60/162 & $5.1{\times}10^{-12}$ & $8.2{\times}10^{-11}$ & $[-19.58,-11.14]$ & 4.82 & 6.93 \\
128 & 2 & solver & $-16.42$ & 59/168 & $2.6{\times}10^{-13}$ & $4.4{\times}10^{-12}$ & $[-20.63,-12.05]$ & 4.67 & 7.08 \\
256 & 0 & solver & $-2.41$ & 110/126 & 0.329 & 1.000 & $[-6.93,+1.51]$ & 4.82 & 7.23 \\
256 & 1 & solver & $-4.67$ & 88/119 & 0.037 & 0.515 & $[-9.04,-0.45]$ & 4.67 & 6.78 \\
256 & 2 & solver & $-1.51$ & 105/115 & 0.544 & 1.000 & $[-5.87,+1.66]$ & 4.82 & 6.93 \\
\addlinespace
\multicolumn{10}{@{}l}{\textit{Config~3 (CodeTAT-QA $\times$ 7B, $n{=}282$)}} \\
2 & 0 & c-SFT & $-3.19$ & 6/15 & 0.078 & 0.940 & $[-6.38,+0.00]$ & 3.90 & 5.32 \\
2 & 1 & c-SFT & $-2.84$ & 9/17 & 0.169 & 1.000 & $[-6.38,-0.35]$ & 4.26 & 5.67 \\
2 & 2 & c-SFT & $-1.77$ & 15/20 & 0.500 & 1.000 & $[-6.03,+1.42]$ & 4.61 & 6.74 \\
64 & 0 & c-SFT & $-2.13$ & 9/15 & 0.307 & 1.000 & $[-5.32,+0.00]$ & 4.26 & 5.67 \\
64 & 1 & c-SFT & $-1.06$ & 10/13 & 0.678 & 1.000 & $[-4.26,+0.71]$ & 3.90 & 5.32 \\
64 & 2 & c-SFT & $-4.26$ & 5/17 & 0.017 & 0.254 & $[-7.45,-1.06]$ & 4.26 & 5.67 \\
\bottomrule
\end{tabular}
\end{table}


Table~\ref{tab:pergroup} reports, for all 17 groups, the McNemar discordant
counts, the exact two-sided $p$, the same $p$ under Holm correction across the
full family of 17, and a paired item bootstrap interval for $\Delta$. These
test the stricter hypothesis, and the reader should keep
Section~\ref{subsec:gotest} in view while reading them.

The honest summary is unflattering. Under Holm correction \textbf{only two
groups are significant}, both of them Config~2 at $m{=}128$
($p_{\text{Holm}} = 8.2{\times}10^{-11}$ and $4.4{\times}10^{-12}$). The
bootstrap interval excludes zero in five groups. Removing the two Config~2
$m{=}128$ groups---the cell where one seed also collapsed---the remaining 15
groups average $-2.91$ points and not one of them is individually significant.
A reader who requires per-group significance should read this paper as
establishing a negative result on one cell and a consistent sign everywhere
else.

That one cell deserves an additional warning against itself. Config~2 at
$m{=}128$ is not merely the cell with the largest effect; it is the cell in
which the RL arm was least healthy. Its two surviving runs have
positive-reward rates of $0.081$ and $0.092$, the lowest of any run in the
study, its third seed collapsed to zero entirely, and its prompt-pool coverage
is $0.19$--$0.22$ epochs. The only two groups that survive multiple-comparison
correction are therefore the two in which under-training is the most plausible
competing explanation. We do not think this invalidates them---continued SFT
and RS-SFT beat GRPO there without the solver's help
(Table~\ref{tab:aggfull}), so the verdict does not rest on the solver---but a
reader should not treat those two $p$-values as the study's strongest evidence.
The strongest evidence is the Go-threshold result of
Section~\ref{subsec:gotest} and the sign pattern below.

The last two columns explain why. Given each group's discordant count, the
smallest $|\Delta|$ that its McNemar test could have declared significant is a
median of $4.82$ points uncorrected and $6.93$ points under Holm. The median
$|\Delta|$ we actually measure is $2.84$ points, and the pre-registered Go
threshold is $+3$. \emph{No single group in this study can resolve a three-point
effect in either direction.}

We want to be exact about what that does and does not concede. It is not a gap
the analysis discovered after the fact: the protocol fixed $+3$ precisely
\emph{because} it put single-seed standard error at $3$--$5$ points
(Section~\ref{sec:method}), so the threshold was knowingly set at the floor of
the design's resolution, and the Go rule was built to compensate---two of three
seeds positive, an \textsc{inconclusive} band over $2$--$5$ points, further
seeds rather than a verdict when signs disagree. The numbers above confirm that
estimate rather than contradicting it: a median MDE of $4.82$ is what an SE of
$3$--$5$ predicts.

What it does concede is symmetric and real. Had GRPO enjoyed a genuine
$+3$-point advantage, this design could not have certified it group by group
either; it would have returned \textsc{inconclusive} and demanded more seeds.
Per-group significance was therefore never the evidence, by construction. The
sign pattern is.

\subsection{What the sign pattern is worth}

All 17 groups are negative. Treating them as independent observations gives a
sign-test $p = 1.5{\times}10^{-5}$, but they are not independent: three seeds
within a cell share a test set, a solver and a task, and the six Config~1 and
Config~3 cells share the test set entirely. We therefore report the same test at
three levels of clustering:

\begin{center}
\begin{tabular}{lrr}
\toprule
unit of independence & units & sign-test $p$ \\
\midrule
group (3 seeds $\times$ 2 budgets $\times$ 3 configs, most optimistic) & 17 & $1.5{\times}10^{-5}$ \\
(config, budget) cell & 6 & $0.031$ \\
configuration & 3 & $0.250$ \\
\bottomrule
\end{tabular}
\end{center}

The defensible reading is the middle row: six independent operating points,
all negative, $p = 0.031$. The headline ``17 negative effect sizes'' is a
description of the data, not seventeen pieces of evidence, and we do not claim
it as such.

\subsection{Is the $\max$ itself producing the sign?}

The sharpest objection to this design is that it commits the error the paper
was written to criticize. $\Delta$ subtracts a $\max$ over four arms, that
$\max$ is taken on the frozen test set, and a maximum over noisy estimates is
upward-biased---so $\Delta$ is biased downward under the null, by an amount
that grows with the number of arms and their noise. At 7B, where the three
learning arms sit within a few points of each other and the whole effect is
$1.06$--$4.26$ points, an analytic max-of-$k$ argument can be made to cover
most of the effect.

That argument assumes the arms are noisy \emph{independently}. They are not:
every arm branches from the same warm start, sees the same items, and fails on
largely the same items, so the per-item correlation between arms is high and
the variance of their maximum is far below the independent case. The question
is empirical, and the per-item vectors settle it.

% 由 analysis/gen_extended_stats.py 生成，勿手改
\begin{table}[t]
\centering
\small
\caption{Winner's-curse audit. \emph{bias} is the bootstrap mean of $\Delta$ with the $\max$ recomputed on each resample minus the same quantity with the arm frozen at the full-sample $\arg\max$. $\Delta_{\text{cf}}$ is cross-fitted: the control arm is selected on a random half of the test set and scored on the other half ($K{=}2{,}000$ splits).}
\label{tab:crossfit}
\begin{tabular}{rrrrrc}
\toprule
$m$ & seed & $\Delta$ (naive) & bias & $\Delta_{\text{cf}}$ & 95\% CI \\
\midrule
\multicolumn{6}{@{}l}{\textit{Config~1 (CodeTAT-QA $\times$ 0.5B, $n{=}282$)}} \\
2 & 0 & $-2.13$ & $+0.000$ & $-2.29$ & $[-8.51,+4.26]$ \\
2 & 1 & $-1.42$ & $+0.000$ & $-1.41$ & $[-7.80,+4.96]$ \\
2 & 2 & $-3.90$ & $+0.000$ & $-3.89$ & $[-9.93,+2.13]$ \\
64 & 0 & $-6.03$ & $+0.000$ & $-6.02$ & $[-12.06,+0.71]$ \\
64 & 1 & $-3.55$ & $-0.002$ & $-3.55$ & $[-9.22,+2.13]$ \\
64 & 2 & $-2.84$ & $+0.000$ & $-2.86$ & $[-9.22,+2.84]$ \\
\addlinespace
\multicolumn{6}{@{}l}{\textit{Config~2 (CodeFinQA $\times$ 1B, $n{=}664$)}} \\
128 & 0 & $-15.36$ & $+0.000$ & $-15.29$ & $[-19.28,-11.14]$ \\
128 & 2 & $-16.42$ & $-0.000$ & $-16.40$ & $[-20.78,-12.05]$ \\
256 & 0 & $-2.41$ & $-0.063$ & $-1.99$ & $[-6.33,+2.71]$ \\
256 & 1 & $-4.67$ & $-0.004$ & $-4.64$ & $[-8.73,-0.30]$ \\
256 & 2 & $-1.51$ & $-0.274$ & $-0.50$ & $[-3.92,+3.31]$ \\
\addlinespace
\multicolumn{6}{@{}l}{\textit{Config~3 (CodeTAT-QA $\times$ 7B, $n{=}282$)}} \\
2 & 0 & $-3.19$ & $-0.002$ & $-3.12$ & $[-6.38,+0.00]$ \\
2 & 1 & $-2.84$ & $-0.364$ & $-2.10$ & $[-4.96,+0.71]$ \\
2 & 2 & $-1.77$ & $-0.172$ & $-1.08$ & $[-4.96,+2.84]$ \\
64 & 0 & $-2.13$ & $-0.200$ & $-1.40$ & $[-4.26,+1.42]$ \\
64 & 1 & $-1.06$ & $-0.482$ & $-0.11$ & $[-2.84,+2.84]$ \\
64 & 2 & $-4.26$ & $-0.003$ & $-4.18$ & $[-7.09,-0.71]$ \\
\bottomrule
\end{tabular}
\end{table}


Table~\ref{tab:crossfit} reports two measurements. The first resamples items
with replacement and computes $\Delta$ twice on each resample---once
recomputing the $\arg\max$, once with the arm frozen at the full-sample
$\arg\max$. The difference is exactly the bias that selecting and scoring on
the same data introduces. It never exceeds $0.472$ points in any group and
averages $-0.091$ points.

The second is a cross-fitted estimator that removes the bias by construction:
the control arm is selected on a random half of the test set and scored on the
other half, averaged over $K = 2{,}000$ splits. Cross-fitted $\Delta$ is
negative in \textbf{17 of 17 groups}, with a mean of $-4.18$ points against the
naive $-4.44$. Selection accounts for a quarter of a point of the four-point
effect.

The reason is visible in Table~\ref{tab:pergroup}: in the 11 groups at 0.5B and
1B the $\arg\max$ is a \emph{deterministic} arm with zero sampling variance
that leads the nearest learning arm by at least $8.5$ points, so no selection
takes place at all. The winner's-curse concern is real only in the six 7B
groups, where three close learning arms compete---and there the measured bias
is at most $0.47$ points against effects of $1.06$ to $4.26$. The objection is
sound in principle and empty in this data.

\subsection{How much of the result is the solver?}

The pre-registered rule aggregates the control group with a $\max$. A reader is
entitled to ask what any weaker aggregation would have said, and what the
result owes to the one arm we built ourselves.

% 由 analysis/gen_extended_stats.py 生成，勿手改
\begin{table}[t]
\centering
\small
\caption{Aggregator sensitivity, summary over the 17 groups. Only the pre-registered $\max$ over the full four-arm control group is negative everywhere; every weaker aggregator reports a positive mean effect.}
\label{tab:aggsummary}
\begin{tabular}{lrrr}
\toprule
control-group aggregator & mean $\Delta$ & range & $\Delta<0$ \\
\midrule
$\max$ over $\mathcal{B}$ (pre-registered) & $-4.44$ & $-16.42$ to $-1.06$ & 17/17 \\
$\max$ over learning arms only & $+5.66$ & $-8.89$ to $+35.11$ & 8/17 \\
second best of $\mathcal{B}$ & $+6.43$ & $-8.89$ to $+35.11$ & 6/17 \\
mean of $\mathcal{B}$ & $+7.97$ & $-6.14$ to $+31.83$ & 2/17 \\
median of $\mathcal{B}$ & $+8.61$ & $-5.80$ to $+40.07$ & 2/17 \\
solver alone & $+1.13$ & $-16.42$ to $+16.31$ & 11/17 \\
RS-SFT alone & $+6.50$ & $-8.89$ to $+35.11$ & 6/17 \\
continued SFT alone & $+9.19$ & $-4.26$ to $+45.04$ & 8/17 \\
warm start $W_m$ alone & $+15.07$ & $+0.71$ to $+48.58$ & 0/17 \\
\bottomrule
\end{tabular}
\end{table}


Table~\ref{tab:aggsummary} answers both. Only the pre-registered $\max$ over
the full four-arm group is negative everywhere. Against the mean of the control
group the study reports $+7.97$ points; against the warm start alone, the
comparison most common in the literature, it reports $+15.07$ points and is
positive in every single group. That spread---from $+15.07$ to $-4.44$ on
identical runs---is the paper's central claim restated as a sensitivity
analysis.

Removing the solver entirely is the most informative row. Against the maximum
over the three learning arms, $\Delta$ averages $+5.66$ points but remains
\emph{negative in 8 of 17 groups}: both Config~2 $m{=}128$ groups and all six
7B groups. In other words, in nearly half the study the verdict does not depend
on the solver at all; continued SFT alone is sufficient to overturn it. The
solver is load-bearing in the 9 remaining groups, all at 0.5B and 1B.
Per-group values are in Appendix~\ref{app:allarms}.

\subsection{The two arms are not solving the same problems}

The discordant counts in Table~\ref{tab:pergroup} carry information the
$\Delta$ summary discards. In Config~1 at $m{=}2$, GRPO is right on 38 items
the solver misses while the solver is right on 44 items GRPO misses: 82 of 282
items are decided differently by the two arms. Their union reaches $78.01\%$
against the better arm's $64.54\%$.

Across all 17 groups the union of GRPO and the best control arm exceeds the
better of the two by $9.02$ points on average (range $1.77$ to $16.57$), and
the union of all five arms averages $75.37\%$ against a best-single-arm average
of $60.67\%$. This is not a deployable method---selecting the right arm per
item requires the answer---but it bounds what the framing costs. Asking which
of a policy-gradient arm and a symbolic arm is stronger, and reporting the
winner, discards the fact that on these tasks they are substantially
complementary. A control-group protocol that reports the $\max$ is still the
right instrument for validating a claimed gain; it is not, on this evidence,
the right instrument for deciding what to build.
